% Angepasste Seminarvorlage
%
\documentclass[12pt,a4paper,fancyheader]{scrartcl}

\usepackage[ngerman]{babel}			% deutsche Namen/Umlaute
\usepackage[utf8]{inputenc}			% Zeichensatzkodierung
\usepackage{csquotes}
\usepackage{fancyhdr}
\usepackage{graphicx}				% Einbinden von Bildern
\usepackage{color}					% Farben wenn es sein muß
\usepackage{hyperref}				% Klickbare Verweise und \autoref{label}
\usepackage{booktabs}				% "Schöne" Tabellen
\usepackage{amsmath}				% Mathematischer Formelsatz AMS
\usepackage{amsfonts}
\usepackage[style=authoryear, backend=biber]{biblatex}
\usepackage{float}     % Für [H] Platzierung
\usepackage{amssymb}
\usepackage{tabularx}
\usepackage{pdflscape}
\usepackage{xparse}
\usepackage[printonlyused]{acronym} % Abkürzungsverzeichnis

% Literatur Daten laden
\addbibresource{libary/literature.bib} % Replace with your bibliography file

% Metadaten Aliasse laden
%
%	Hier werden Titel, Bearbeiter und das Datum eingetragen
%
\newcommand\papertopic{\LaTeX-Vorlage für Seminararbeiten an der FOM}
\newcommand\paperauthors{Max Mustermann, Erika Musterfrau}
\newcommand\paperdate{\today} % \today für das aktuelle Datum, sonst als Text eintragen
\newcommand\paperabstract{
  Hier steht eine fünf bis zehn Zeilen lange Zusammenfassung der Arbeit.
  \blindtext
}

\newcommand\coursename{Modul: Seminararbeiten und andere Möglichkeiten Studenten zu foltern}
\newcommand\coursesemester{WiSe 2025}
\newcommand\courseinstitution{FOM Hochschule für Oekonomie \& Management}


% Template Befehle laden
% Das ist ein 100 Wörter umfassender Blindtext, um das Aussehen von längeren Texten zu testen.
\newcommand{\blindtext}{\textit{Lorem ipsum dolor sit amet, consetetur sadipscing elitr, sed diam nonumy eirmod tempor invidunt ut labore et dolore magna aliquyam erat, sed diam voluptua. At vero eos et accusam et justo duo dolores et ea rebum. Stet clita kasd gubergren, no sea takimata sanctus est Lorem ipsum dolor sit amet. Lorem ipsum dolor sit amet, consetetur sadipscing elitr, sed diam nonumy eirmod tempor invidunt ut labore et dolore magna aliquyam erat, sed diam voluptua. At vero eos et accusam et justo duo dolores et ea rebum. Stet clita kasd gubergren, no sea takimata sanctus est Lorem ipsum dolor sit amet.}}

% citation commands
\newcommand{\dircite}[2][]{\footcite[][#1]{#2}}
\newcommand{\indircite}[2][]{\footcite[Vgl.][#1]{#2}}
\newcommand{\dircitegraph}[2][]{
\\{\small \textcite{#2}\IfNoValueF{#1}{\unskip,~#1}}
}
\newcommand{\indircitegraph}[2][]{
\\{\small In~Anlehnung~an:~\textcite{#2}\IfNoValueF{#1}{\unskip,~#1}}%
}


% Abkürzungsverzeichnis laden
\newcommand{\abkuerzungen}{
    \begin{acronym}[FOM] % längste Abkürzung in eckigen Klammern
        \acro{FOM}{Fachhochschule für Oekonomie \& Management}
        \acro{IT}{Informationstechnologie}
    \end{acronym}
}

\hypersetup{ % Thema und Author in die Meta-Daten der PDF
pdftitle={Seminararbeit zum Thema: \papertopic} 
pdfauthor={\paperauthors}
}

\begin{document}
\begin{titlepage}
    \centering

    % Titel der Arbeit
    {\huge\bfseries \papertopic \par}
    \vspace{1cm}

    % Autor und Datum
    \begin{minipage}{\textwidth}
        \centering
        {\Large \textsl{\paperauthors}\par}
        \vspace{1cm}
        \hfill {\Large \paperdate}
    \end{minipage}
    \vspace{1.5cm}

    % Logo
    \includegraphics[width=0.45\textwidth]{images/FOM.jpg}\par\vspace{6pt}

    % Kursinfos
    {\large \textsc{\coursename} \par}
    {\coursesemester \, · \, \courseinstitution \par}
    \vspace{1cm}

    \vfill

    % Abstract Bereich
    \begin{minipage}{0.85\textwidth}
        \textbf{Abstract:}\par\vspace{0.25\baselineskip}
        \small
        \setlength{\parindent}{1.5em}
        \paperabstract
    \end{minipage}

\end{titlepage}
%\thispagestyle{empty} % lässt die Seitennummer auf der Titelseite verschwinden


\pagenumbering{Roman}
\cleardoublepage

\tableofcontents			% Inhaltsverzeichnis ... nicht für kurze Dokumente!
\cleardoublepage

\listoffigures				% Abbildungsverzeichnis ... nicht für kurze Dokumente
\cleardoublepage

\pagenumbering{arabic}

% Text der Arbeit
%
%   List an tex Dateien, die Inhalt der Arbeit enthalten.
%   canary.tex enthält eine Anlaitung und probiert die meisten features aus.
%

%
%	z.B. Text des ersten Autors
%

\section{Einführung}
Eine Untergliederung des Dokuments in Abschnitte und Unterabschnitte
geschieht mit den Kommandos \verb|\section{...}| und \verb|\subsection{...}|,
in Absätze mit \verb|\par| oder einfach einer Leerzeile im \LaTeX-Quelltext.

\paragraph{Betitelter Absatz.}
Ein sehr kleiner Abschnitt (meistens nur ein Absatz)
mit Titeltext, der nicht vom Textrumpf abgesetzt ist, wird durch das Kommando
\verb|\paragraph{titel}| erzeugt.

Bitte seien Sie mit dem Kommando \verb?\verb|text|? vorsichtig! Die Anweisung bewirkt eine exakte Darstellung des Textes darin.

%
%	z.B. Text des zweiten Autors
%

\section{Weitere Kommandos}

\subsection{Mathematische Formeln}\label{sec:formeln}

Mathematische Formeln werden mittels \verb|\(...\)|
in den Fließtext eingebaut
--- zum Beispiel \( E=mc^2 \) und:
sei \(V\) ein Vektorraum über \(\mathbb{R}\)
und \(\mathcal{M}\) eine Indexmenge
---
oder aber mittels \verb|\[...\]|
abgesetzt und zentriert dargestellt:
\[
    \pmb{x} = \sqrt[3]{\frac{a^2-b^2}{a^2+b^2}}
    ~~~~~ \text{versus} ~~~~~
    \boldsymbol{x} = \sqrt[3]{\frac{a^2-b^2}{a^2+b^2}}
\]


\subsection{Silbentrennung}\label{sec:silbentrennung}
Vertrauen Sie bitte nie einer automatischen Silbentrennung (auch nicht der von Microsoft Word \& Co.). In folgendem Test-Absatz ist das Wort "`Spracherkennung"' falsch getrennt.

Testzeile Testzeile Testzeile Testzeile Testzeile Testzeile Testzeile Testzeile Te Spracherkennung.

Sie können LaTeX die richtige Trennung mit \verb|\hyphenation{...}| mitteilen. Man tut das üblicherweise noch vor \verb|\begin{document}|.

\hyphenation{Sprach-er-ken-nung}	% hier jetzt zu Demonstrationszwecken im Text
Testzeile Testzeile Testzeile Testzeile Testzeile Testzeile Testzeile Testzeile Te Spracherkennung.


\subsection{Literaturzitate}
\label{sec:literaturzitate}

Diese Vorlage basiert auf einer Latex Vorlage von Claudius Stern\indircite{Stern2018b} mit meinen eigenen Anpassungen.
Konkret wurden extra Zitierbefehle angelegt, um angepasst an die Anforderungen aus dem Kurs Wissenschaftliches Arbeiten zitieren zu können.

Das Literaturverwaltungsprogramm JabRef\indircite{JabRef2014} ist für viele Plattformen verfügbar und unterstützt bei der Literaturrecherche. Es ist prädestiniert dazu, mit {\LaTeX} in Kombination mit {Bib\TeX} zusammenzuarbeiten.

So kann in dieser Vorlage mit \verb|\dircite{Bibtexkey}| eine Quelle direkt Zitiert werden.
Optional kann auch ein Seitenhinweis mit \verb|\dircite[S. 42]{Buch}| wie folgt beigefügt werden.
Indirekte Zitate können analog mit \verb|\indircite[page info]{Bibtexkey}| erstellt werden.

\subsection{Einbinden von Bildern}

Sie können mit der Anweisung \verb|\includegraphics{datei}| eine Grafikdatei einbinden. Diese kann im PDF-, JPG- oder PNG-Format vorliegen. Grafiken werden üblicherweise in die float-Umgebung \verb|figure| gekapselt. Jede eingebundene Grafik muss vom Text aus referenziert werden. Die textuelle Referenz hat \textit{vor} der Grafik zu erfolgen.

\begin{figure}[tb]
    \centering
    \includegraphics{images/FOM.jpg}
    \dircitegraph{FOM}
    \caption{Logo der FOM}
    \label{fig:fomlogo}
\end{figure}

Auch andere Grafikformate werden unterstützt.
Grafiken finden sich nicht notwendigerweise direkt am Einfüge-Ort (siehe \autoref{fig:fomlogo}).
Das Textsatzsystem richtet es so ein, dass es gut aussieht.
Zitate von Graphiken sind möglich über \verb|\dircitegraph[]{}| und \verb|\indircitegraph[]{}|.

\subsection{Querverweise}\label{sec:querverweise}
Querverweise können sie mit dem \verb|\label{key}| befehl direkt nach einer section Anweisung oder als Teil einer Graphik erstellen.
Auf diese können sie dann anschließend \verb|\autoref{key}|, oder wenn sie nur die Nummer brauchen mit \verb|\ref{key}| verweisen.

Näheres zu Formeln können sie Beispielsweise in \autoref{sec:formeln} lesen, während in \ref{sec:silbentrennung} Silbentrennung erklärt wird.

\subsection{Abkürzungen}
Abkürzung können nach einem Eintrag in der Datei \verb|abbreviation.tex| mit den Befehlen \verb|\ac{XYZ}| (für einen kleinen ersten Buchstaben) und \verb|\Ac{XYZ}| (für einen großen ersten Buchstaben) in den Text eingefügt werden.

Dieses Template ist für Seminararbeiten an der \ac{FOM} gedacht.
Es wurde für einen \ac{IT} Studiengang der \ac{FOM} angepasst.

\subsection{Erstellen von Tabellen}

Das Volk hat gesprochen. Siehe \autoref{tab:tabular}. Auch hier kommt ein float zum Einsatz, jedoch mit dem Positionierungs-Hinweis \verb|[h]|. Jedes float sollte mit einem Label versehen werden, und es sollte im Text darauf verwiesen werden, da sich die Position ändern kann.

\begin{table}[h]
    \centering
    \begin{tabular} {|llc||r|}
        \hline
        Name       & Rang    & Fraktion & Stimmenanteil \\
        \hline
        Mobutu     & General & CDU      & 57\%          \\
        Tsvangirai & Oberst  & CSU      & 63\%          \\
        \hline
    \end{tabular}
    \caption {Bundestagswahl in Simbabwe}
    \label{tab:tabular}
\end{table}

\autoref{tab:booktab} verwendet keine vertikalen Linien und entspricht dem üblicherweise in Büchern verwenden Stil.
Derartige Tabellen sind deutlich ansehnlicher.

\begin{table}[h]
    \centering
    \begin{tabular} {llcr}
        \toprule
        Name       & Rang    & Fraktion & Stimmenanteil \\
        \midrule
        Mobutu     & General & CDU      & 57\%          \\
        Tsvangirai & Oberst  & CSU      & 63\%          \\
        \bottomrule
    \end{tabular}
    \caption {Bundestagswahl in Simbabwe}
    \label{tab:booktab}
\end{table}

\subsection{Listen und Aufzählungen}

Listen und Aufzählungen werden in einer Umgebung angelegt (umschlossen von \verb|begin| und \verb|end|.).
\verb|\begin{itemize}| leitet eine Liste ein und \verb|\begin{enumerate}| eine Aufzählung.
Die Einträge werden jeweils mit \verb|\item| begonnen.
Folgend zwei Beispiele.

Itemize:
\begin{itemize}
    \item Test
    \item Test
    \item Test
\end{itemize}

Enumerate:
\begin{enumerate}
    \item Test
    \item Test
    \item Test
\end{enumerate}
 % enthält einen Referenztext und stellt einige oft genutzte Befehle vor.

% Das Literaturverzeichnis
\cleardoublepage
\addcontentsline{toc}{section}{Literaturverzeichnis}
\printbibliography[title={Literaturverzeichnis}]

% Das Abkürzungsverzeichnis
\cleardoublepage
\section*{Abkürzungsverzeichnis}
\addcontentsline{toc}{section}{Abkürzungsverzeichnis} % optional für TOC
\abkuerzungen
\end{document}
